\section{Proving a Pin Without Trusting Your Eyes}
\label{sec:gpio-proving}

The blink is the first program every embedded developer writes, and it is the first one that lies. A
blink that does not blink has four different causes that look identical: the code never ran, the pin
is wrong, the polarity is wrong, or the LED is not connected. This section is about telling them
apart with instruments that are already there.

\subsection{The Case That Made This Necessary}

This project had a blink that flashed successfully and did not blink. The write verified, the
read-back matched, the board was silent. Every step of the toolchain reported success:

\begin{minted}{text}
cargo build            Finished `release` profile [optimized] target(s)
python3 scripts/elf-load-range.py   load_min=0x08006000 load_max=0x08006813
probe-rs download      Finished in 0.41s
probe-rs verify        Finished in 0.28s
\end{minted}

Four green steps and no light. The only way forward was to stop asking the board to show and start
asking the register to tell.

\subsection{Reading the Latch Over SWD}

The probe can read any address while the program runs. The output latch is one word, and it toggles
or it does not:

\begin{minted}{bash}
# halt, read DOUT twice, resume
openocd -c "init; halt" -c "mdw 0x5003C040 1" -c "sleep 500" -c "mdw 0x5003C040 1" -c "resume; shutdown"

# 0x5003c040: 00000080      <- bit 7 high
# 0x5003c040: 00000000      <- bit 7 low, 500 ms later
\end{minted}

That pair of readings is a proof of the \emph{program}: the latch is moving at the rate the code
intended. What it is not is a proof of the \emph{board}: a latch that toggles and an LED that stays
dark means the pin is not the LED's pin, or the LED is active-low and the code drives it backwards.

The second reading of \codeinline{PORTA\_MODEL} closes the other half:

\begin{minted}{text}
PORTA_MODEL = 0x00000040   # mode 4 (push-pull) in the field for pin 7
\end{minted}

Latch moving, mode correct, and still no light: that combination is now a board question, not a
software one — which is exactly the division of labour a measurement is for.

\subsection{Three Instruments, Three Questions}

\begin{table}[htbp]
    \centering
    \footnotesize
    \begin{tabular}{@{}p{0.26\textwidth}p{0.34\textwidth}p{0.34\textwidth}@{}}
        \toprule
        \textbf{Instrument} & \textbf{Question it answers} & \textbf{How it is read} \\
        \midrule
        The LED & is the whole chain alive? & eyes, and only after the two below \\
        \codeinline{DOUT}/\codeinline{MODEL} over SWD & is the program right? & \codeinline{mdw} on the port \\
        A mailbox in RAM & did the code reach this point? & \codeinline{mdw} on a magic word \\
        \bottomrule
    \end{tabular}
    \caption{Order matters: the LED answers last, because it answers slowly and ambiguously.}
    \label{tab:gpio-instruments}
\end{table}

\subsection{The Mailbox, Which Is the Real Lesson}

The instrument that made the rest of this project possible is not a pin at all. It is a small block
of RAM — a magic word, a sequence number, a status and a message — that a program writes before it
does anything it might not survive:

\begin{minted}{rust}
// 0x2003F000: magic "MBOX", seq, status (0 info / 1 pass / 2 fail), len, message bytes
mailbox::write(Status::Pass, b"cap=0x20000000 4000@400Hz real=clean");
\end{minted}

Read with one \codeinline{mdw} over SWD, it answers questions that no LED can: did the classifier
load its calibration, how many samples did the ring hold, which of the two I2C aliases answered. It
costs 300 bytes of RAM and it replaced an entire class of guesswork — which is why every example in
this book reports through it rather than through a serial port that may not exist.
